\section{Glossary}\label{sec:glossary}

This glossary is a reader's guide to the technical vocabulary used in the Catalogue that follows.  It says, in language that is as plain as possible, what each term is doing in this paper; for each substantive formula it glosses the symbols, says why the formula has that form, and gives a small worked example with real numbers; and it points to implementation details only where the code fixes the meaning.  One notational warning up front: the paper reserves the letter $F$ for the per-tick variational free energy and for the Bayesian-model-reduction evidence change $\Delta F$; the cascade objective is deliberately called a \emph{score}, because it is an engineering selection quantity, not a free energy.

\paragraph{Active Inference Framework.} Active inference is a theory of adaptive behaviour: an agent keeps a model of the world, compares that model with what it senses, and chooses actions expected to make future evidence less surprising or more useful \cite{FRISTON2016862,friston2019freeenergyprincipleparticular}.  In this paper the phrase is used operationally.  The War Machine is not a brain model or a replacement coding agent, but a harness that observes, updates beliefs, scores possible actions, acts when gated, and observes again, when wrapped around existing agentic coding systems.\footnote{The scheduled loop is implemented in \srcref{futon2/scripts/wm_scheduled_run.clj}; task scoring is in \srcref{futon2/src/futon2/aif/efe.clj}; the live-loop baseline is recorded in \srcref{futon2/holes/wm-baseline.md}.}

\paragraph{Generative model.} A generative model is the system's story about how hidden causes produce observable evidence.  In a coding-agent setting, the hidden causes are operational states such as progress, risk, blockage, uncertainty, or pattern fit; the observations are things such as test results, trace records, diffs, and operator feedback \cite{FRISTON2016862,neacsu2024structure}.\footnote{Belief channels and preference scoring are implemented across \srcref{futon2/src/futon2/aif/belief.clj}, \srcref{futon2/src/futon2/aif/preferences.clj}, and \srcref{futon2/src/futon2/aif/efe.clj}.}

\paragraph{Belief state $\mu$.} The Greek letter $\mu$ names the system's current best estimate of its situation.  Instead of asking the language model to reread the entire conversation every time, the system stores a compact map of operational hypotheses: what goal is active, what is blocked, what looks uncertain, and what kind of pattern may fit.  The continuous tutorial form tracks a mean and an exponential-moving-average variance with a sensor-noise floor. The production War Machine does not store that pair: it maintains a categorical posterior over entity statuses, applies declared likelihood events, and derives observation-channel means and variances from that posterior.\footnote{R1/R8 are described in \srcref{futon2/holes/aif-r1-r16-pattern-map.md}; the categorical filter, updater, and derived channel predictions are implemented in \srcref{futon2/src/futon2/aif/belief.clj}.}
% QUESTION: How does the "belief state" relate to memory?  Memory seems to be things that were believed to be true at one point, and also worth remembering because they are likely to be true (or at least relevant) in the future.  Since there is a structural analogy I think it would be worth talking about how this does or does not relate to the model.

\paragraph{Observation vector $o$.} An observation vector is a standardized summary of what happened.  The system cannot compare ``tests failed,'' ``operator objected,'' and ``diff got larger'' unless these events are normalized into named channels.  That is why the catalogue emphasizes typed observations rather than raw prose.  As measured on 2026-08-31, the live observation vector has fourteen channels, eight of which have an R3a likelihood model (a $14\to8$ projection); these counts are derived in the implementation tests from the channel collections rather than maintained as a second list in prose.\footnote{The authorities are \texttt{observation-channels} in \srcref{futon2/src/futon2/aif/observation.clj} and \texttt{channels-with-likelihood} in \srcref{futon2/src/futon2/aif/belief.clj}; their derived counts are checked in \srcref{futon2/test/futon2/aif/observation_test.clj} and \srcref{futon2/test/futon2/aif/belief_test.clj}. Trace fields are guarded in \srcref{futon2/src/futon2/aif/trace.clj}.}
% QUESTION: This also relates to the question of "compare over time".  In an agentic coding system, "raw prose" is often enough if the prose and observations fit within the agent's context window.  But the WM is set up to run over long periods of time, in different sessions, with different agents that do not share a context window.  Although the observation vector has to do with what happened at a given moment, it seems to relate, also, to establishing a shared context over time.

\paragraph{Prediction error $\varepsilon$.} Prediction error is the gap between what the system expected and what it observed.  If the system expected progress but tests fail, the error is large.  In the update rule $\mu \leftarrow \mu + \alpha\,\Pi\,\varepsilon$, read the arrow as ``replace the old belief by a corrected belief'': the correction is the error $\varepsilon$, scaled by a learning rate $\alpha$ and by precision $\Pi$, the trust placed in that error signal \cite{FRISTON2016862}.  This is one gradient step of predictive coding --- exactly the correction a Kalman filter makes, with $\alpha\,\Pi$ playing the role of the Kalman gain.  A worked example: $\mu=0.50$, $o=0.90$, $\alpha=0.3$, $\Pi=2$ give $\varepsilon=0.40$ and $\mu \leftarrow 0.50+0.3\cdot2\cdot0.40=0.74$; a high precision makes the step larger, which is why untrusted channels are given small $\Pi$.\footnote{Predictive-coding style updates are specified in the R3 pattern and tied to precision channels in \srcref{futon2/src/futon2/aif/belief.clj}.}

\paragraph{Precision $\Pi$.} Precision is confidence expressed as a weight.  A high-precision signal is trusted more; a low-precision signal is discounted.  This matters because software evidence is uneven: a passing unit test, a stale heuristic, and an LLM explanation should not all count equally.\footnote{Per-channel evidence precision is the R7 pattern; selection gain $g$ is the R14 pattern and is logged through the War Machine trace path.}

\paragraph{Variational free energy $F$.} Variational free energy is a present-tense mismatch score: how badly the current belief explains the current observation.  In $F = \tfrac{1}{2}\,\mathrm{mean}_k(\Pi_k\,\varepsilon_k^2)$\eqanchor{F}, the subscript $k$ means ``for each evidence channel'' --- tests, logs, traces, operator feedback, and so on.  Squared error makes misses positive and penalizes large misses; precision weights trusted channels more heavily.  It is not a moral or thermodynamic claim here; it is an operational alarm and diagnostic number.  If $F$ jumps, the system's model and evidence are out of alignment \cite{FRISTON2016862,friston2019freeenergyprincipleparticular}.  The form is the Laplace/Gaussian collapse of variational free energy to a precision-weighted sum of squared errors, with the conventional Gaussian $\tfrac12$.  Note that the same $\varepsilon$ drives the belief update above: $F$ is the \emph{size} of the mismatch, the update is the \emph{move} that shrinks it --- and where $F$ scores the present fit, $G$ (next entry) scores imagined futures.  A worked example: $\varepsilon=[0.2,\,1.0]$ with $\Pi=[4,\,1]$ gives $F=\tfrac12\,\mathrm{mean}(0.16,\ 1.0)=0.29$.\footnote{Trace persistence is guarded in \srcref{futon2/src/futon2/aif/trace.clj}.}

\paragraph{Expected free energy $G$.} Expected free energy is a future-tense action score: if the system takes this action, or follows this policy, how much risk, ambiguity, and useful learning should it expect?  Written out, $G_{\mathrm{efe}}(a) = D_{KL}[Q(o\mid a)\Vert C] + \mathbb{E}\,H[P(o\mid s)]$\eqanchor{Ga}, in nats: the first term is risk, the divergence of predicted outcomes from preferences $C$; the second is ambiguity, the expected observation entropy (both are glossed separately below).  The implementation persists this two-term core apart from its controller augmentation --- posterior spread, urgency, intrinsic credit, feasibility masks, and coverage bonuses are reported beside it, never hidden inside it.  So risk $8$ and ambiguity $6$ give $G_{\mathrm{efe}}=14$, and a controller total of $11$ means the explicitly reported bonuses contributed $-3$.  The paper uses two grains: $G(a)$ scores candidate tasks, while $S(\pi)$ scores short tactical cascades of patterns for the chosen task \cite{FRISTON2016862,neacsu2024structure}.\footnote{Task-level scoring is in \srcref{futon2/src/futon2/aif/efe.clj}; tactical construction scoring uses \srcref{futon2/src/futon2/aif/fold_eval.clj} and the fold escrow path.}

\paragraph{Predictive outcome distribution $Q(o\mid\pi)$.} $Q(o\mid\pi)$ is the probability distribution over outcomes the model predicts if policy $\pi$ is followed; the task-level formula above writes the same idea as $Q(o\mid a)$ for a candidate action $a$.

\paragraph{Preference distribution $C$.} $C$ is the probability distribution over outcomes the agent prefers, against which its predictive outcome distribution is compared inside the KL-divergence risk term.  This is the canonical active-inference preference distribution, not the implementation's separately named vector of pragmatic costs.

\paragraph{Risk.} Risk is the expected cost of ending up in an unwanted state.  The implementation uses the canonical active-inference form: a KL divergence between predicted outcomes and preferred outcomes.\footnote{\srcref{futon2/src/futon2/aif/efe.clj} supports \texttt{:risk-mode :kl}; \srcref{futon2/test/futon2/aif/efe_risk_mode_test.clj} pins the default and the escape hatch.}

\paragraph{Ambiguity.} Ambiguity is uncertainty in the mapping from state to outcome.  The notation $P(o\mid s)$\eqanchor{ambiguity} is read ``the probability of observing $o$ if the world is in state $s$.''  Its entropy is high when that state could produce many different observations, and low when the observation is predictable.  Expected entropy means averaging that uncertainty over the states the system may encounter.  Reducing ambiguity means choosing actions that make future evidence more informative \cite{FRISTON2016862}.  For a binary outcome, $P(o\mid s)=[0.5,0.5]$ gives entropy $1$ bit --- maximum ambiguity --- while $[0.9,0.1]$ gives about $0.47$ bit.\footnote{\srcref{futon2/src/futon2/aif/efe.clj} supports \texttt{:ambiguity-mode :gaussian-entropy}; tests are in \srcref{futon2/test/futon2/aif/efe_test.clj}.}

\paragraph{Observation model $A$.} The observation model $A$ is the agent's theory of its sensors: for each hidden state $s$, it specifies the probability distribution over observations $P(o \mid s)$. In the production War Machine belief path, every event update flows through this declared model --- when the agent observes $o$, it revises its posterior by multiplying the predicted prior by the likelihood column $A[o \mid \cdot]$. The implementation declares $A$ as an explicit, column-normalised $7 \times 7$ matrix over the entity-status vocabulary, with three entry classes: diagonal entries (an observation is the strongest evidence for its own status), lifecycle-adjacent entries (an observation is mildly compatible with a neighbouring status), and contradictory entries (an observation is evidence \emph{against} an opposed status). The contradictory class lets the model say ``this observation counts against that status,'' not merely ``for this one.'' The transition model $B$ is the identity, so the prediction step preserves the prior posterior, and the initial prior $D$ is uniform. The filter computes the prediction--update cycle $q^- = Bq$, $q \propto A(o \mid \cdot)^{\kappa(w)} q^-$, where $\kappa(w) = \log_2(1+w)$ is a tempered-likelihood exponent controlled by observation weight. Each trace stamps the selected likelihood mode and content hashes for A, B, and D; a provenance-visible diagonal comparison mode supports controlled ablation. A simulation over 50 synthetic entities with varied event types confirms that the off-diagonal structure is behaviourally active: it redistributes posterior mass at $50{\times}$ the rate of the real trace corpus (KL mean 0.050 vs 0.001 nats), and correctly suppresses contradictory statuses (observing \texttt{strengthened} reduces \texttt{falsified} mass by 33\% relative to the diagonal comparison). This warrants R3's green, model-relative status: the filter is exactly derived from its declared A/B/D model. It does not establish that those hand-set entries describe the world; calibration against exogenous evidence from adjudicated handoffs, build outcomes, and operator review remains open.\footnote{The declared observation model, transition model, validators, categorical filter, model manifest, and production seam are in \srcref{futon2/src/futon2/aif/belief.clj} and \srcref{futon2/scripts/futon2/report/war_machine.clj}; tests are in \srcref{futon2/test/futon2/aif/belief_test.clj} and \srcref{futon2/test/futon2/aif/a_matrix_live_wiring_test.clj}; the simulation is \srcref{futon2/scripts/aif/a_matrix_simulation.bb}; design and findings are in \srcref{futon2/holes/missions/M-aif-a-matrix-faithfulness.md} and \srcref{futon2/holes/labs/M-aif-faithfulness/a-matrix-simulation-findings.md}.}

\paragraph{Model uncertainty and EIG.} Expected information gain is an expectation of posterior uncertainty reduction. The structure learner does not compute that expectation: it exposes aggregate posterior spread, $U_{\mathrm{model}} = \sum_c \mathrm{sd}(A_c)$\eqanchor{eig-stddev} over the concepts' Dirichlet posteriors, as a \emph{model-uncertainty bonus}. The bonus may steer exploration inside the engineering controller, but it is deliberately excluded from \texttt{G-efe}: it does not predict an observation, nor the posterior-entropy reduction that observation would cause. A policy-conditioned EIG term would require predicted prior and posterior entropies under each policy.  The boundary is worth stating twice\eqanchor{efe-offline}: online $G_{\mathrm{efe}}$ ranks predicted outcomes, offline Bayesian Model Reduction (below) compares model evidence over stored observations, and posterior spread may suggest which reduction to test --- but none of that turns either quantity into EIG.\footnote{\texttt{:model-uncertainty-bonus} is composed in \srcref{futon2/src/futon2/aif/efe.clj}; the posterior-spread helpers are in \srcref{futon2/src/futon2/aif/a4a.clj} and \srcref{futon2/src/futon2/aif/a4a_substrate.clj}.}

\paragraph{Softmax and controller calibration.} Softmax turns scores into a probability distribution: $p_i = e^{-G_i/\tau}\big/\sum_j e^{-G_j/\tau}$\eqanchor{softmax}.  Because lower scores are better, softmax is applied to the negated scores, so low-cost actions become high-probability picks; the temperature $\tau$ sets how decisively --- low $\tau$ sharpens toward the best option, high $\tau$ flattens and explores.  A worked example: $G=[1,2,4]$ at $\tau=1$ gives $p\approx[70\%,\,26\%,\,4\%]$, and halving $\tau$ to $0.5$ pushes the best option toward $88\%$. The effective temperature is set through the outcome-feedback selection gain $g$, as $\tau_{\mathrm{eff}}=1/g$; score-spread normalisation is disabled. This is an engineering calibration control, not an inferred variational policy precision.  At the selection seam the learned habit prior enters unscaled by temperature --- the choice is over $-G/\tau_{\mathrm{eff}} + \ln E(\pi)$ --- and each decision records, in its decision explanation, the per-term contributions and which of the two terms governed it.\footnote{Action ranking and calibration are discussed in R14; the scheduled loop records the selected action, score breakdown, and decision explanation in the trace.}

\paragraph{Pattern language / cascade.} A design pattern is a named response to a recurring problem, written so other people can recognize and reuse it.  This paper makes patterns operational: the same pattern text can guide a human reader and serve as an action candidate or warrant for an agent \cite{alexander1977a,alexander1999a,gamma1994a,corneli2015patterns,corneli2021patterns,corneli2022patterns,corneli2026patternsnewgenerationai}.  In this paper, \emph{cascade} is the operational name for a pattern language in use: a staged composition of patterns fitted to the selected mission.  It says which warrants are being combined, and in what operational order, before any work is enacted.  The harness is not merely choosing an isolated next command; it is proposing a small pattern-language construction that can be priced, checked, folded, and only then enacted.\footnote{The active-inference catalogue keys each pattern to an R-number; historical pattern material is tracked in the prior PLoP papers cited in \srcref{p4ng/main.bib}. \srcref{futon2/scripts/futon2/report/cascade_lane.clj} builds the cascade-policy: \texttt{cascade-policy-for} constructs it, and \texttt{cascade-lane} records \texttt{:shown [pattern-ids...]}, \texttt{:cascade-score}, and \texttt{:policy-rollout-score}.}

\paragraph{Control states $U$ and the policy vocabulary.} In a discrete active-inference model, $U$ is the set of actions, also called \emph{control states}; a policy $\pi$ is an allowable sequence drawn from those controls, and the $E$ vector below places a prior over the already-allowable policies.  This paper works one level higher: a design pattern is treated as an elementary or temporally extended \emph{control schema}, while a cascade is the policy composed from those schemas.  The correspondence is functional rather than literal --- AIF does not know about design-pattern prose --- but it locates the important difference between two kinds of novelty. Cascade construction searches new combinations inside a fixed vocabulary $U$; pattern authoring proposes $U' = U \cup \{u^*\}$ and must also propose how the new control changes expected states, $B' = B \cup \{B^{u^*}\}$, before a forward model can score policies containing it.  The induced horizon-$H$ policy space therefore grows, $\Pi_H(U) \subset \Pi_H(U')$.  Saying that a new pattern enters at a ``uniform prior'' is an engineering shorthand: canonical $E$ has one entry per complete allowable policy, whereas the implementation's pattern-level proposal/reliability prior is a separate quantity that may eventually induce priors over cascades containing the pattern.  R17$^{\prime\prime\prime}$ names this policy-vocabulary expansion and its residual explicitly \cite{friston2019freeenergyprincipleparticular,smith2020structure}.

\paragraph{Policy prior $E$ and habit.} Canonical $E(\pi)$ is the prior probability of an already-allowable policy before the current observations are used to judge it.  It is commonly read as habit: repeated selection can increase a policy's prior probability without establishing that the policy is good.  At the selection seam this prior has the characteristic unscaled seat
\[
  Q(\pi) \;\propto\; \exp\!\left(\ln E(\pi)-G(\pi)/\tau\right),
\]
so temperature or policy precision modulates the expected-free-energy term, not the habit prior.  The War Machine implements this form exactly but not yet at the canonical policy grain.  It folds selected scheduler actions into stable categories $k(a)=(\text{type},\text{target or target-class})$ and uses a symmetric Dirichlet posterior predictive over the currently feasible menu,
\[
  \widehat E(a_i) = \frac{n_{k(a_i)}+\alpha}{\sum_{j\in\Pi_{\mathrm{feasible}}}(n_{k(a_j)}+\alpha)},
  \qquad \alpha=1,
\]
with no recency decay; duplicate identities divide their category mass.  The resulting $\ln\widehat E$ enters selection unscaled and is persisted with its sufficient statistics and decision explanation.  This is a real learned frequency prior, and live traces show it can govern a choice even when $G$ prefers another candidate.  The residual is grain, not algebra: its categories are mission/action-class identities rather than complete pattern cascades, and the tactical cascade lane does not yet consume it.  Consequently \emph{pattern enters at the uniform prior}, \emph{cascade prior}, and this scheduler-level $\widehat E$ must not be treated as synonyms.\footnote{The pure fold and posterior predictive are in \srcref{futon2/src/futon2/aif/habit_prior.clj}; \srcref{futon2/src/futon2/aif/policy.clj} implements $\ln\widehat E-G/\tau$; live wiring and persistence are in \srcref{futon2/scripts/futon2/report/war_machine.clj}.  The measured case in which the learned prior overruled the $G$ ranking is reconstructed in \srcref{futon2/holes/TN-wm-rank109-explained.md}.}

\paragraph{Policy $\pi$.} In this paper, a policy $\pi$ is the active-inference name for a pattern language/cascade when that composition is being scored.  Its units are design patterns and their foldable warrants, not generic commands.  It is tempting to read the composition as purely temporal --- do this pattern, then this one, then this one --- but that is the impoverished, tree-shaped reading.  After Alexander's \emph{a city is not a tree}, and its stack-scale analogue \emph{a proof is not a tree}, the cascade is a \emph{semilattice} of patterns: they compose both by sequential dependency (\texttt{BV.seq}, the ``then'') and by cross-cutting co-application (\texttt{BV.copar}, the overlapping meets where a single pattern belongs to several sub-constructions at once).  Empirically the overlap axis dominates --- a coverage-selected cascade carries a handful of sequential edges against many more co-application meets, and some carry no ``then'' at all.  In $S(\pi)=\sum_t \rho^t s(s_t)$\eqanchor{Gpi}, the summation sign means ``add this over the cascade steps.''  The small $s(s_t)$ is the cost of the state reached at step $t$; the large $S(\pi)$ is the total engineering rollout score of the cascade-policy.  The factor $\rho^t$ discounts later steps, so nearer consequences count more --- the same discounted-return idea as in reinforcement learning, letting a distant payoff matter without swamping the immediate situation.  (For $s=[3,2,4]$ and $\rho=0.9$: $3 + 0.9\cdot2 + 0.81\cdot4 = 8.04$.)  This temporal sum scores only the sequential axis; the co-application structure is what the \emph{fold} turns into a wiring diagram, and what the cascade's \emph{wholeness} (coherence $\times$ intensity, after Alexander) measures.  The distinction earns its keep twice over: a pattern that looks poor as an isolated first move may be good as the opening of a longer construction, and a cascade flattened to a linear list folds to an \emph{empty} wiring where the same cascade folded \emph{as a semilattice} yields a full construction (one box per pattern, its meets as \texttt{copar}).  Cascade size is then not a fixed budget but a model choice, $\arg\max_\pi S_{\mathrm{cascade}}(\pi)$ with $S_{\mathrm{cascade}}=\mathrm{coverage\ reward}-\lambda\cdot\mathrm{prior\ cost}$\eqanchor{cascade-F} (coverage reward $5$ against prior cost $2$ at $\lambda=1$ scores $3$): an everything-cascade scores negative --- in information-theoretic terms, selecting nothing accomplishes nothing.\footnote{R13 addresses this one-step-vs-policy-horizon issue: the gap is between scoring only the next action and scoring the pattern-language cascade/policy it opens. \srcref{futon2/src/futon2/aif/close_loop.clj} turns a cascade-lane entry into an act-gate with \texttt{:cascade-score} and \texttt{:coverage-score-delta}; \srcref{futon2/src/futon2/aif/enact.clj} carries \texttt{:shown} into enactment. Fold/cascade checking uses \srcref{futon2/src/futon2/aif/fold_llm.clj}, \srcref{futon2/src/futon2/aif/fold_escrow.clj}, and \srcref{futon2/scripts/fold_author.clj}.  The cascade's semilattice (descent $=$ \texttt{BV.seq}, co-application $=$ \texttt{BV.copar}) is constructed in \srcref{futon3a/holes/labs/M-memes-arrows/cascade_construct.py} (\texttt{construct\_cascade}, \texttt{chosen\_semi\_lattice}) over the pattern phylogeny \srcref{futon6/data/pattern-phylogeny-edges.json}, which emits an engineering cascade score (coverage reward minus prior cost) and Alexander wholeness at each cascade size.}

\paragraph{Aliveness $L = T\cdot H$.} Salingaros's measure of a structure's ``life'' multiplies a temperature-like intensity $T$ (detail, curvature, colour) by a harmony factor $H$ (coherence, internal symmetry); $T=0.8$ and $H=0.6$ give $L=0.48$ (see the Background for provenance and the caveat that this is not claimed to be a free-energy functional).  Its role here is concrete: the discharge-trained proposer credits the realised $L$ of a \emph{built} cascade back to the patterns it used, so aliveness enters the loop as a learned reward rather than an aesthetic aside.

\paragraph{Embedding space.} An embedding turns text into a point in a large numerical space.  Nearby points often have related meanings.  The paper uses this as a generator of hypotheses: nearby patterns may be candidates for composition or comparison, but nearness alone is not proof.\footnote{The constellation data used by A4a comes from \srcref{futon3c/holes/excursions/pipeline-semilattice-clusters.edn}.}

\paragraph{Bayesian Model Reduction.} Bayesian Model Reduction asks whether a simpler model explains the accumulated evidence well enough to replace a more complex one.  In $A' = A + a' - a$\eqanchor{bmr-post}, $a$ is the old prior, $A$ is the old posterior after seeing evidence, $a'$ is the reduced model's prior, and $A'$ is the reduced model's posterior.  These are vectors of smoothed counts, so the arithmetic is element-by-element.  The formula reuses the evidence already accumulated, but expresses it under the reduced model's assumptions; indeed the form is forced by the counts themselves: since $A = a + (\text{counts})$, substituting the reduced prior gives $A' = a' + (\text{counts}) = A + a' - a$, with no new data collected.  A worked example: $a=[1,1]$, $A=[10,4]$, and reduced prior $a'=[1,0.01]$ give $A'=[10,\,3.01]$.  In this paper BMR is used for structure learning: can two concepts be merged, or can a redundant parameter be pruned, without collecting new data \cite{neacsu2024structure}?\footnote{\srcref{futon2/src/futon2/aif/bmr.clj} implements the pure BMR kernel; \srcref{futon2/test/futon2/aif/bmr_test.clj} pins the sign convention and accept/reject behaviour.}

\paragraph{Dirichlet concentration parameters.} A Dirichlet distribution is a probability model over probability tables.  Its concentration parameters behave like smoothed counts: larger counts mean stronger evidence.  The paper uses this because pattern concepts and mission outcomes can be represented as count tables \cite{neacsu2024structure}.\footnote{\texttt{dirichlet-moments} in \srcref{futon2/src/futon2/aif/bmr.clj} returns per-component means and standard deviations; A4a aggregates those standard deviations as endpoint-count units.}

\paragraph{Log multivariate beta and $\Delta F$.} The log multivariate beta function is part of the closed-form evidence calculation for Dirichlet models.  Readers do not need to compute it by hand.  It is the accounting device that lets the system compare a full model with a reduced model analytically: $\Delta F = \ln B(A) + \ln B(a') - \ln B(a) - \ln B(A')$, where $B(\cdot)$ is the multivariate beta function --- the normaliser of a Dirichlet --- and the four terms are the log-normalisers of the old and reduced priors and posteriors \cite{neacsu2024structure}.  This $\Delta F$ is the \emph{reduction} evidence change, distinct from the per-tick variational $F$ above.  The logarithm is used because the underlying probabilities can be tiny; logs keep the arithmetic stable and turn products into sums.\footnote{\texttt{log-multivariate-beta} and \texttt{bayesian-model-reduction} are in \srcref{futon2/src/futon2/aif/bmr.clj}; tests cover identity, redundant-parameter acceptance, and informative-parameter rejection.}

\paragraph{Bayes factor threshold.} The threshold $\Delta F \le -3$\eqanchor{bmr-threshold} is a rule of thumb for requiring substantial evidence before simplifying the model.  In this paper's sign convention, a sufficiently negative $\Delta F$ means the reduced model wins.  Three nats corresponds to about $\exp(3)\approx20$, so the threshold says, informally, ``only merge when the evidence is about twenty-to-one in favour of the simpler model.''  So $\Delta F=-3.5$ passes (about thirty-three to one) while $\Delta F=-2.0$ does not.  It prevents the system from merging concepts just because they look superficially similar.\footnote{The threshold is encoded in \srcref{futon2/src/futon2/aif/bmr.clj}; A4a's count-BMR over-merging result is recorded as a negative finding in \srcref{futon2/holes/labs/slush-demo/findings/}.}

\paragraph{GFlowNet ``slush''.} A GFlowNet is a sampler that learns to generate many high-reward objects rather than only the single best one.  In the paper, the slush proposes diverse pattern compositions --- cascades --- for the selected mission, sampling a cascade $S$ for mission $m$ with probability $P(S \mid m) \propto \exp(\beta\,\hat R(S\mid m))$, where $\hat R$ is the learned reward and the inverse temperature $\beta$ is kept deliberately low so the distribution stays broad.  Diversity is the point: a greedy (high-$\beta$) proposer yields a monoculture, and the loop learns fastest from varied proposals.  A worked example: for $\hat R=[2,1]$, $\beta=0.5$ gives $P\approx[0.62,\,0.38]$, while $\beta=1$ already sharpens to $\approx[0.73,\,0.27]$.  In offline tests the slush beats a greedy baseline on diversity and coverage, while its corpus-global information term does not steer above uniform.  The trainer uses a per-mission $\log Z$ and recovers an exactly enumerated Gibbs distribution (total-variation $0.0125$).  Its reward is learned from adjudicated fold outcomes because pattern-text coverage cannot distinguish a well-written plan from one that works: on the adjudicated record, a flown-and-killed cascade out-covered most true dischargers. The loop therefore uses the slush as a low-temperature diversity supplier rather than treating text coverage as reward.\footnote{Offline slush artifacts and the review/spec line live in \srcref{futon2/holes/labs/slush-demo/} (see \srcref{futon2/holes/TN-gflownets-fable-review.md} at \srcref{futon2/holes/} and \srcref{futon2/holes/labs/slush-demo/SPEC-full-loop-gfn.md}); the trusted trainer is \srcref{futon2/holes/labs/slush-demo/slice2/gfn_core.py}; the producer-side live-loop lookup is \texttt{slush-candidates} in \srcref{futon2/src/futon2/aif/a4a_substrate.clj}.}

\paragraph{Fold.} A fold is a checked construction plan.  It turns a sequence of patterns into a small wiring diagram: boxes are steps, wires show how outputs feed inputs, terminals name the desired result, and policy-holes mark honest gaps.  The point is to make the agent's proposed action checkable instead of merely plausible.  In everyday terms, it is like a recipe or a circuit board: each step has a defined input and output, and the whole plan can be inspected before anything runs.\footnote{\texttt{fold-prompt} is in \srcref{futon2/src/futon2/aif/fold_llm.clj}; escrow validation is in \srcref{futon2/src/futon2/aif/fold_escrow.clj}; authoring support is in \srcref{futon2/scripts/fold_author.clj}.}

\paragraph{Act-gate.} The act-gate is the last check before anything runs: act if and only if $S_{\mathrm{cascade}} > 0 \,\wedge\, \Delta S_{\mathrm{coverage}} < 0$\eqanchor{gate} --- the cascade's coverage reward exceeds its prior cost, \emph{and} the fold improves the lower-is-better coverage score.  These are engineering control quantities, not $F$ or $G$; if either leg is absent, the gate abstains rather than guesses.  So $S_{\mathrm{cascade}}=+1.2$ with $\Delta S_{\mathrm{coverage}}=-0.8$ passes, while a positive coverage-score delta fails.\footnote{\srcref{futon2/src/futon2/aif/close_loop.clj} emits \texttt{:cascade-score} and \texttt{:coverage-score-delta}.}

\paragraph{EDN} EDN is a data notation from the Clojure ecosystem, similar in purpose to JSON but with richer data types.  Using EDN here means the agent's construction can be parsed and checked as data, not treated as prose.\footnote{Fold deposits in \srcref{futon6/data/fold-turns/} are EDN records checked by \srcref{futon2/src/futon2/aif/fold_escrow.clj}.}

\paragraph{Substrate and Drawbridge.} The paper uses \emph{substrate} at two layers.  Substrate-1 is the direct working world: files, Git commits, logs, test results, and other filesystem-level artifacts that ordinary tools can read and write directly.  Substrate-2 is the graph-database model of that world in XTDB: entities, hyperedges, capabilities, discharges, stars, and other typed records derived from or referring back to substrate-1.  Drawbridge is the guarded access layer to substrate-2, not a general filesystem wrapper.  This matters because an agent should not be allowed to certify its own internal story; it must change or cite a shared store and then re-observe that store through an independent path.  ``No self-certification'' means the agent cannot simply declare success: it has to produce an artifact or graph record that another process, gate, test, or later observation can independently inspect.\footnote{Drawbridge helpers for substrate-2 are in \srcref{futon2/src/futon2/aif/actuator_a3.clj}; guarded star/candidate writes are in \srcref{futon2/src/futon2/aif/a4a_substrate.clj}. Filesystem artifacts such as fold deposits live directly on substrate-1, for example \srcref{futon6/data/fold-turns/}.}

\paragraph{No self-certification.} No self-certification means the system cannot count evidence it manufactured as if it were independent proof.  This is the core auditability rule: a claim must be backed by a witness outside the part of the system making the claim.  Calibration inherits the same rule in two layers: $L_1$ compares the model's predicted $G$ against its own realised $G$ --- cheap, run every cycle, but self-referential, so it may never be reported as value evidence --- while $L_2$ compares predictions against outcomes the model did \emph{not} produce, and only $L_2$ certifies that the model's value is right.  A model can pass $L_1$ forever while being worthless at predicting real outcomes; only a referent the model had no hand in escapes ``prior $=$ value.''\footnote{The birth-tagging of evidence records is implemented in \srcref{futon0/scripts/futon0/futonzero/rollout_ledger.clj} and \srcref{futon0/scripts/futon0/futonzero/reward_red_team.clj}, with tag-discipline tests in \srcref{futon3c/test/futon3c/aif/flight_record_test.clj}; corroboration of review \emph{execution} from the job-event stream is in \srcref{futon2/src/futon2/aif/full_loop_runner.clj}. Trace provenance is guarded in \srcref{futon2/src/futon2/aif/trace.clj}. \emph{Caution on numbering:} this paper's R-numbers are its own, and do \emph{not} align with the R-numbers of the \srcref{futon2} AIF-completeness contract, where R9 denotes named validation properties and R6 denotes softmax action selection. Earlier drafts of this footnote cited \srcref{futon2/test/futon2/aif/r9_named_validation_test.clj} on the strength of the matching numeral; that file tests the \srcref{futon2} criterion, not this pattern.}

\paragraph{Demonstration Foundry and have--want arrows.} A have--want arrow says: here is a capability or situation we have, and here is the state we want.  Demonstration Foundry extends the scan loop to public repositories by looking for compositions between our arrows and others' arrows.  A pull request is only one outcome; adopting an existing tool or rejecting a false match can also be a valid result.  This connects the active-inference loop to commons and peer-production settings \cite{benkler2006wealth,wenger1998communities}.\footnote{Mission design is in \srcref{futon7/holes/M-demonstration-foundry.md}; meme-arrow background is in \srcref{futon3a/README-memes-and-arrows.md}; daily scan code begins in \srcref{futon7/src/f7/daily.clj}.}

\paragraph{Strategic mission selection.} The slow side of the hierarchy has a second grain problem.  The live War Machine enriches each mission with an additive engineering value --- centrality, fit to a declared strategy cascade, and phase doability --- before task-level EFE selection.  This quantity is not $E_S$: centrality is structural relevance, strategy fit is context-conditioned proposal support, and doability is a provisional predictor of progress.  Nor does multiplying the three numbers repair the theory.  The implementation explicitly rejected that product because a mission on the strategy spine can have measured centrality zero, which would annihilate the very strategic evidence being introduced; more fundamentally, multiplying three hand-set $[0,1]$ scores does not make a generative model.  This failure is diagnostic: the scalars are surrogates for outer-loop structure that the implementation has not yet represented.  A principled strategic layer instead distinguishes (i) reason-bearing policy support, including hard executability and operator gates; (ii) predicted mission outcomes and their uncertainty, which belong in the forward model and $G_S$; (iii) context-conditioned proposal potentials; and (iv) the separate frequency habit $E_S(\pi_S)$.  The p4ng catalogue can organize this landscape as a small control ontology.  The War Machine may write concrete, provenance-bearing episodic memories in the Zaif style, attached by typed hyperedges to missions and R-patterns; an outer cascade of control patterns then acts directly as an ordered retrieval and candidate-construction policy.  This reuses cascade composition, holes, and explanations at the slower grain instead of compressing them into ``centrality.''  A dedicated embedding over missions plus these high-level control patterns remains a possible exploratory proposal mechanism, but is neither foundational nor required: direct pattern-conditioned recall works without vector search, and typed hypergraph evidence disposes.  The resulting R15 hierarchy selects a strategic mission-policy by $Q(\pi_S)\propto E_S(\pi_S)\exp[-G_S(\pi_S)/\tau_S]$, then selects a tactical pattern cascade conditionally by $Q(\pi_T\mid\pi_S)\propto E_T(\pi_T\mid\pi_S)\exp[-G_T(\pi_T\mid\pi_S)/\tau_T]$, with separate identities, evidence, and learned state at the two grains.\footnote{The current additive model is specified in \srcref{futon2/holes/M-wm-three-factor-mission-value.md}.  The typed strategic landscape, factor-role audit, memory route, and staged replacement are in \srcref{futon2/holes/missions/M-wm-strategic-mission-selection.md}; its dark schema kernel is \srcref{futon2/src/futon2/aif/mission_control_graph.clj}.}

\paragraph{Clicks, attempts, and cohorts (measurement vocabulary).} The system itself runs one loop: observe $\to$ select $\to$ construct $\to$ dispatch $\to$ build $\to$ review $\to$ adjudicate.  A \emph{click} is a single trigger of that loop.  An \emph{attempt} is the complete forward pass a click produces --- each attempt begins with a fresh selection against the current state of the field, so attempts are numbered in click order but carry no retry relationship to one another; work that fails leaves a typed repair obligation \emph{in the field}, which a later attempt's selection may pick up as ordinary first-class work (the loop re-encounters its debt as terrain; it never rolls back).  The only backward motion is \emph{within} a single attempt: one bounded build-cure retry, and one bounded revise-and-resubmit round after reviewer findings.  A \emph{cohort} is not part of the mechanism at all --- it is the measurement discipline laid over it: a preregistered window (``count the next $N$ attempts under the stated semantic epoch'') whose stopping rule, outcome taxonomy, and no-relabelling terms are committed before any attempt runs, so that every failure stays in the denominator and results are never pooled across changes to the machine's semantics.  The empirics in Section~\ref{sec:empirics} report attempts; cohorts are why those reports can be trusted.\footnote{Cohort preregistrations live in \srcref{futon2/holes/labs/M-aif-full-loop-46/cohort.edn}; the append-only attempt records and stopping rule are implemented in \srcref{futon2/src/futon2/aif/full_loop_cohort.clj}.}

\paragraph{Revision (2026-08-31; cancellation boundary).} The cancellation outcome was added to the runner after cohort 46 preregistered its outcome taxonomy.  A cancelled attempt remains durably visible, but it may not be pooled under that original taxonomy: cancellation begins a new semantic stratum unless and until a cohort preregisters it explicitly.  The implementation's global acceptance of \texttt{:cancelled} is therefore not itself evidence that cohort 46 anticipated that outcome.

\paragraph{A shared experimental substrate.} The mathematics learner and the War Machine are consequently application domains of one control architecture, not separate memory theories.  Zaif runners can act as ``nomadic'' workers: they traverse problem territories and leave typed, revisable episodic traces in the same bitemporal hypergraph that supports WM reflection.  Domain, subject, author, pattern, mission, provenance, and witness endpoints preserve the boundary between the territories even when storage and retrieval machinery are shared.  This gives a cost-ordered validation ladder: pure projection tests, frequent mathematics runs, dark WM decisions, and finally expensive live WM clicks.  Mathematics can cheaply test whether realization, durable writing, pattern-conditioned recall, cascade composition, use receipts, correction, and supersession work.  It cannot certify WM-specific preferences, transition predictions, safety, or mission value; those claims still require independent observations from the WM domain.  The transferable object is the learning mechanism and its calibration discipline, not the truth of domain-specific outcomes.
